\documentclass[../../main.tex]{subfiles}

\graphicspath{{../../image/}{../../figure/}}

\begin{document}

\section{相对熵与经典解的稳定性}\label{section:ch05-相对熵与经典解的稳定性}

本节旨在证明：凸熵的存在保证了初值问题的经典解连续依赖于初值，即使在更广泛的可容许有界弱解类中亦然。

\begin{theorem}\label{theorem:ch05-相对熵稳定性}
假设守恒律方程组\eqref{align:ch05-cauchy-system}装备有熵--熵通量对 $(\eta,Q)$，其中 $\mathrm{D}^{2}\eta(U)$ 在 $\overline{\mathscr{B}_{\rho}}$ 上正定。设 $\bar{U}$ 是\eqref{align:ch05-cauchy-system}在 $[0,T)$ 上的经典解，取值于 $\overline{\mathscr{B}_{\rho}}$，初值为 $\bar{U}_{0}$。设 $U$ 是\eqref{align:ch05-cauchy-system}在 $[0,T)$ 上的任一弱解，取值于 $\overline{\mathscr{B}_{\rho}}$，满足熵可容许性条件\eqref{align:ch04-熵积分0}，
且初值为 $U_{0}$。则对任意 $r>0$ 与 $t\in[0,T)$，成立
\begin{align}
\int_{|x|<r}|U(x,t)-\bar{U}(x,t)|^{2}\,\mathrm{d}x\leqslant ae^{bt}\int_{|x|<r+st}|U_{0}(x)-\bar{U}_{0}(x)|^{2}\,\mathrm{d}x,\label{align:ch05-relative-stability}
\end{align}
其中正常数 $s,a$ 仅依赖于 $G,\eta,Q$ 及其在 $\overline{\mathscr{B}_{\rho}}$ 上导数的界，而 $b$ 还依赖于 $\bar{U}$ 的 Lipschitz 常数。特别地，$\bar{U}$ 是以 $\bar{U}_{0}$ 为初值、取值于 $\overline{\mathscr{B}_{\rho}}$ 的唯一可容许弱解。
\end{theorem}

\begin{proof}
在 $\overline{\mathscr{B}_{\rho}}\times\overline{\mathscr{B}_{\rho}}$ 上定义函数
\begin{align}
h(U,\bar{U})=\eta(U)-\eta(\bar{U})-\mathrm{D}\eta(\bar{U})[U-\bar{U}],\label{align:ch05-h-def}
\end{align}
\begin{align}
Y_{\alpha}(U,\bar{U})=Q_{\alpha}(U)-Q_{\alpha}(\bar{U})-\mathrm{D}\eta(\bar{U})[G_{\alpha}(U)-G_{\alpha}(\bar{U})],\label{align:ch05-Yalpha-def}
\end{align}
\begin{align}
Z_{\alpha}(U,\bar{U})=A(\bar{U})\{G_{\alpha}(U)-G_{\alpha}(\bar{U})-\mathrm{D}G_{\alpha}(\bar{U})[U-\bar{U}]\},\label{align:ch05-Zalpha-def}
\end{align}
它们关于 $U-\bar{U}$ 均为二阶量（回忆\eqref{align:ch04-熵流关系}与\eqref{align:ch05-A-def}）。
因此，由于 $\mathrm{D}^{2}\eta(U)$ 在 $\overline{\mathscr{B}_{\rho}}$ 上正定，存在正常数 $s$ 使得
\begin{align}
|Y(U,\bar{U})|\leqslant sh(U,\bar{U}).\label{align:ch05-Y-bound}
\end{align}

固定任意非负、Lipschitz 连续、在 $\mathbb{R}^{m}\times[0,T)$ 上具紧支集的试验函数 $\psi$，沿两个解 $U(x,t)$、$\bar{U}(x,t)$ 计算 $h,Y,Z$。由于 $U$ 满足不等式\eqref{align:ch04-熵积分0}，而 $\bar{U}$ 作为经典解恒等满足\eqref{align:ch04-熵积分0}为等式，我们推得
\begin{align}
\int_{0}^{T}\int_{\mathbb{R}^{m}}[\partial_{t}\psi h(U,\bar{U})+\sum_{\alpha=1}^{m}\partial_{\alpha}\psi Y_{\alpha}(U,\bar{U})]\,\mathrm{d}x\,\mathrm{d}t+\int_{\mathbb{R}^{m}}\psi(x,0)h(U_{0}(x),\bar{U}_{0}(x))\,\mathrm{d}x
&\geqslant-\int_{0}^{T}\int_{\mathbb{R}^{m}}\{\partial_{t}\psi\mathrm{D}\eta(\bar{U})[U-\bar{U}]+\sum_{\alpha=1}^{m}\partial_{\alpha}\psi\mathrm{D}\eta(\bar{U})[G_{\alpha}(U)-G_{\alpha}(\bar{U})]\}\,\mathrm{d}x\,\mathrm{d}t\nonumber\\
&\quad-\int_{\mathbb{R}^{m}}\psi(x,0)\mathrm{D}\eta(\bar{U}_{0}(x))[U_{0}(x)-\bar{U}_{0}(x)]\,\mathrm{d}x.\label{align:ch05-relative-ineq1}
\end{align}

接下来对两个解 $U$ 与 $\bar{U}$ 分别写出\eqref{align:ch04-弱解0}，
以 Lipschitz 连续向量场 $\psi\mathrm{D}\eta(\bar{U})$ 为试验函数 $\Phi$，得
\begin{align}
\int_{0}^{T}\int_{\mathbb{R}^{m}}\{\partial_{t}[\psi\mathrm{D}\eta(\bar{U})][U-\bar{U}]+\sum_{\alpha=1}^{m}\partial_{\alpha}[\psi\mathrm{D}\eta(\bar{U})][G_{\alpha}(U)-G_{\alpha}(\bar{U})]\}\,\mathrm{d}x\,\mathrm{d}t+\int_{\mathbb{R}^{m}}\psi(x,0)\mathrm{D}\eta(\bar{U}_{0}(x))[U_{0}(x)-\bar{U}_{0}(x)]\,\mathrm{d}x=0.\label{align:ch05-relative-id1}
\end{align}
由于 $\bar{U}$ 是\eqref{align:ch05-cauchy-system}的经典解，并由\eqref{align:ch05-A-def}、\eqref{align:ch05-Jalpha-def}，
\begin{align}
\partial_{t}\mathrm{D}\eta(\bar{U})=\partial_{t}\bar{U}^{\top}A(\bar{U})=-\sum_{\alpha=1}^{m}\partial_{\alpha}\bar{U}^{\top}J_{\alpha}(\bar{U})^{\top}=-\sum_{\alpha=1}^{m}\partial_{\alpha}\bar{U}^{\top}A(\bar{U})\mathrm{D}G_{\alpha}(\bar{U}),\label{align:ch05-Deta-t}
\end{align}
从而，回忆\eqref{align:ch05-Zalpha-def}，
\begin{align}
\partial_{t}\mathrm{D}\eta(\bar{U})[U-\bar{U}]+\sum_{\alpha=1}^{m}\partial_{\alpha}\mathrm{D}\eta(\bar{U})[G_{\alpha}(U)-G_{\alpha}(\bar{U})]=\sum_{\alpha=1}^{m}\partial_{\alpha}\bar{U}^{\top}Z_{\alpha}(U,\bar{U}).\label{align:ch05-Z-identity}
\end{align}

合并\eqref{align:ch05-relative-ineq1}、\eqref{align:ch05-relative-id1}与\eqref{align:ch05-Z-identity}，得
\begin{align}
\int_{0}^{T}\int_{\mathbb{R}^{m}}[\partial_{t}\psi h(U,\bar{U})+\sum_{\alpha=1}^{m}\partial_{\alpha}\psi Y_{\alpha}(U,\bar{U})]\,\mathrm{d}x\,\mathrm{d}t+\int_{\mathbb{R}^{m}}\psi(x,0)h(U_{0}(x),\bar{U}_{0}(x))\,\mathrm{d}x
\geqslant\int_{0}^{T}\int_{\mathbb{R}^{m}}\psi\sum_{\alpha=1}^{m}\partial_{\alpha}\bar{U}^{\top}Z_{\alpha}(U,\bar{U})\,\mathrm{d}x\,\mathrm{d}t.\label{align:ch05-relative-ineq2}
\end{align}

现固定 $t\in(0,T)$ 与 $r>0$。对任意 $\sigma\in(0,t]$ 与充分小的正数 $\varepsilon$，对试验函数 $\psi(x,\tau)=\chi(x,\tau)\omega(\tau)$ 写出\eqref{align:ch05-relative-ineq2}，其中
\begin{align}
\omega(\tau)=\begin{cases}1,&0\leqslant\tau<\sigma,\\ \varepsilon^{-1}(\sigma-\tau)+1,&\sigma\leqslant\tau<\sigma+\varepsilon,\\ 0,&\sigma+\varepsilon\leqslant\tau<\infty,\end{cases}\label{align:ch05-cutoff-t}
\end{align}
\begin{align}
\chi(x,\tau)=\begin{cases}1,&|x|-r-s(\sigma-\tau)<0,\\ \varepsilon^{-1}[r+s(t-\tau)-|x|]+1,&0\leqslant|x|-r-s(t-\tau)<\varepsilon,\\ 0,&|x|-r-s(t-\tau)\geqslant\varepsilon,\end{cases}\label{align:ch05-cutoff-x}
\end{align}
其中 $s$ 是\eqref{align:ch05-Y-bound}中出现的常数。计算给出
\begin{align}
\frac{1}{\varepsilon}\int_{\sigma}^{\sigma+\varepsilon}\int_{|x|<r+s(t-\sigma)}h(U(x,\tau),\bar{U}(x,\tau))\,\mathrm{d}x\,\mathrm{d}\tau
&\leqslant\int_{|x|<r+st}h(U_{0}(x),\bar{U}_{0}(x))\,\mathrm{d}x\nonumber\\
&\quad-\frac{1}{\varepsilon}\int_{0}^{\sigma}\int_{r+s(t-\tau)<|x|<r+s(t-\tau)+\varepsilon}\left[sh(U,\bar{U})+\frac{Y(U,\bar{U})x}{|x|}\right]\,\mathrm{d}x\,\mathrm{d}\tau\nonumber\\
&\quad-\int_{0}^{\sigma}\int_{|x|<r+s(t-\tau)}\sum_{\alpha=1}^{m}\partial_{\alpha}\bar{U}^{\top}Z_{\alpha}(U,\bar{U})\,\mathrm{d}x\,\mathrm{d}\tau+O(\varepsilon).\label{align:ch05-cutoff-integral}
\end{align}
令 $\varepsilon\downarrow0$。由\eqref{align:ch05-Y-bound}，\eqref{align:ch05-cutoff-integral}右端第二个积分非负。故
\begin{align}
\int_{|x|<r+s(t-\sigma)}h(U(x,\sigma),\bar{U}(x,\sigma))\,\mathrm{d}x
&\leqslant\int_{|x|<r+st}h(U_{0}(x),\bar{U}_{0}(x))\,\mathrm{d}x\nonumber\\
&\quad-\int_{0}^{\sigma}\int_{|x|<r+s(t-\tau)}\sum_{\alpha=1}^{m}\partial_{\alpha}\bar{U}^{\top}Z_{\alpha}(U,\bar{U})\,\mathrm{d}x\,\mathrm{d}\tau,\label{align:ch05-after-cutoff}
\end{align}
对 $L^{\infty}$ 弱$^{*}$ 连续性 $\eta(U(\cdot,\tau))$ 在 $(0,t)$ 中的所有点 $\sigma$ 成立。如前所述，$h(U,\bar{U})$ 与 $Z_{\alpha}(U,\bar{U})$ 关于 $U-\bar{U}$ 均为二阶量，此外 $h(U,\bar{U})$ 因 $\eta$ 凸而正定。于是令
\begin{align}
u(\tau)=\int_{|x|<r+s(t-\tau)}|U(x,\tau)-\bar{U}(x,\tau)|^{2}\,\mathrm{d}x,\label{align:ch05-u-def}
\end{align}
\eqref{align:ch05-after-cutoff}蕴涵
\begin{align}
u(\sigma)\leqslant au(0)+b\int_{0}^{\sigma}u(\tau)\,\mathrm{d}\tau,\label{align:ch05-Gronwall}
\end{align}
对几乎所有 $\sigma\in(0,t)$ 成立。由于 $u(\cdot)$ 弱下半连续，\eqref{align:ch05-Gronwall}对所有 $\sigma\in[0,t]$ 成立。由 Gronwall 不等式得 $u(t)\leqslant au(0)e^{bt}$，即\eqref{align:ch05-relative-stability}。注意 $a$ 与 $s$ 仅依赖于 $G,\eta,Q$ 及其在 $\overline{\mathscr{B}_{\rho}}$ 上导数的界，而 $b$ 还依赖于 $\bar{U}$ 的 Lipschitz 常数。证明完成。
\end{proof}

值得注意的是：只要经典解存在，单个凸熵的熵不等式就能筛出初值问题唯一的解。然而正如我们将看到的，当经典解不存在时，单个熵不等式一般不再足以挑出某一特定弱解。特别地，如\cref{section:ch04-Euler方程}所见，
Euler 方程组\eqref{align:ch03-等熵Euler}
在特殊构造的初值下，允许无穷多个满足关于熵 $\rho\varepsilon(\rho)+\frac{1}{2}\rho|v|^{2}$ 以等式成立的熵可容许性条件\eqref{align:ch04-熵积分0}的弱解。
弱解唯一性问题错综复杂，将是后续各章讨论的主要课题。

由\eqref{align:ch05-h-def}与\eqref{align:ch05-Yalpha-def}定义的 $U$ 的函数 $h(U,\bar{U})$ 与 $Y(U,\bar{U})$ 通常称为相对于状态 $\bar{U}$ 的\emph{相对熵}（relative entropy）及相应的\emph{相对熵通量}（relative entropy flux）。

\begin{remark}\label{remark:ch05-相对熵正定}
在\cref{theorem:ch05-相对熵稳定性}的证明中，仅需 $h(U,\bar{U})$ 对经典解值域中的所有 $\bar{U}$ 正定即可。即使 $\eta$ 不凸，当经典解特殊时这也可能成立；例如，当 $\bar{U}$ 是 $\eta$ 的强极小值的常状态。
\end{remark}

\end{document}
